\documentclass[12pt]{article}
\usepackage[utf8]{inputenc}
\usepackage[T1]{fontenc}
\usepackage[margin=1in]{geometry}
\usepackage{setspace}
\usepackage{xcolor}
\usepackage[colorlinks=true, urlcolor=blue!60!black, linkcolor=blue!60!black]{hyperref}
\setlength{\parindent}{2em}
\setlength{\parskip}{1.25em}
\renewcommand{\baselinestretch}{1.0}
\begin{document}
\begin{center}
{\LARGE \textit{Wisconsin v. Yoder}, 406 U.S. 205 (1972)\par}
\vspace{0.8em}
{\large Case Brief --- Four-Page Edition\par}
\end{center}
\vspace{1em}

\section{Memory Jogger}

Applying strict scrutiny under the Free Exercise Clause, the Court held that Wisconsin could not criminally punish Old Order Amish and Conservative Amish Mennonite parents for refusing to send their children to formal high school after the eighth grade, because the State's interest in one or two more years of schooling did not outweigh the grave threat to the Amish religious way of life (\href{https://supreme.justia.com/cases/federal/us/406/205/}{406 U.S. 205}).

\section{Detailed Case Facts}

Jonas Yoder and Wallace Miller (Old Order Amish) and Adin Yutzy (Conservative Amish Mennonite) of Green County, Wisconsin, declined to enroll their children, aged 14 and 15, in high school after eighth grade. Wisconsin required attendance until age 16, on pain of fines or up to three months' imprisonment. The State Superintendent rejected a compromise modeled on the ``Pennsylvania plan'' of part-time Amish vocational schooling.

The State stipulated that the parents' beliefs were sincere. Uncontradicted experts, Dr.\ John Hostetler and Dr.\ Donald Erickson, testified that Amish faith, rooted in Anabaptism and the command ``be not conformed to this world,'' pervades daily life through the \emph{Ordnung}. The Amish accepted elementary schooling for basic literacy but opposed high school because it removed adolescents from the community before adult baptism and exposed them to worldly values; Erickson called Amish ``learning through doing'' ``ideal.'' Only Frieda Yoder testified, confirming her own religious objection.

\section{Procedural History}

The parents were convicted in Green County Court and fined \$5 each; the court found a burden on sincere beliefs but upheld the law as ``reasonable and constitutional.'' The Circuit Court affirmed, the Wisconsin Supreme Court reversed under the Free Exercise Clause, and the U.S. Supreme Court affirmed on May 15, 1972.

\section{Judicial Votes}

Affirmed, \textbf{6--1}, with seven Justices participating.
\begin{itemize}
\item \textbf{Majority:} Chief Justice Burger, joined by Justices Brennan, Stewart, White, Marshall, and Blackmun.
\item \textbf{Concurrences:} Justice Stewart (joined by Brennan); Justice White (joined by Brennan and Stewart).
\item \textbf{Dissenting in part:} Justice Douglas.
\item \textbf{Not participating:} Justices Powell and Rehnquist.
\end{itemize}

\section{Holding}

The First and Fourteenth Amendments bar Wisconsin from compelling Amish parents to send their children to formal high school until 16. Only state ``interests of the highest order and those not otherwise served'' can overbalance legitimate Free Exercise claims. Given the severe burden on a centuries-old religious way of life, the Amish acceptance of eight years of schooling and an effective vocational alternative, and no shown harm to the children or society, Wisconsin failed to show with particularity how a one- or two-year exemption would impair its goals. \textit{Parens patriae} authority under \textit{Prince v. Massachusetts} does not reach this record.

\section{Analysis of Opinions}

\textbf{Majority (Burger).} Though ``[p]roviding public schools ranks at the very apex of the function of a State,'' that interest is subject to balancing, and \textit{Pierce} linked free exercise to parents' interest in their children's religious upbringing. A secular way of life like Thoreau's at Walden ``does not rise to the demands of the Religion Clauses,'' but the Amish way of life is ``one of deep religious conviction, shared by an organized group, and intimately related to daily living.'' The Court rejected the belief--action dichotomy, since belief and action ``cannot be neatly confined in logic-tight compartments,'' and the claim that facial neutrality suffices. Examining the marginal effect of an exemption, it found the Amish law-abiding and self-sufficient, Jefferson's ``sturdy yeoman,'' and the fear that departing youth would burden society ``highly speculative.'' \textit{Prince} involved threats to safety or health; \textit{Pierce} ``stands as a charter of the rights of parents to direct the religious upbringing of their children.'' No parent--child conflict was before the Court, and the Amish had made a ``convincing showing, one that probably few other religious groups or sects could make.''

\textbf{Stewart (concurring).} The case concerned only criminal punishment of parents, not children's right to attend high school.

\textbf{White (concurring).} A ``close'' case; it would be ``a very different case'' if a religion forbade all schooling, but the deviation here was ``relatively slight.''

\textbf{Douglas (dissenting in part).} Mature children are ``persons'' whose views must be heard; ``[i]t is the future of the student, not the future of the parents, that is imperiled.'' The Amish record of order is irrelevant, and excluding Thoreau's conscience conflicts with \textit{Seeger} and \textit{Welsh}.

\textbf{Later developments.} \textit{Smith} (1990) preserved \textit{Yoder} as a ``hybrid situation,'' and \textit{Mahmoud v. Taylor} (2025) revitalized it: ``[w]e have never confined \emph{Yoder} to its facts.''

\section{Examples: Future Applications}

\begin{enumerate}
\item \textbf{Same-Side (Exemption Required): Attendance Age 18.} A State raises the age to 18 and prosecutes Amish parents whose teens learn farming at home. The \textit{Yoder} burden is the same, and the State's interests are already served.
\item \textbf{Same-Side (Strict Scrutiny): Normative Instruction.} Parents seek opt-outs from a unit designed to change their children's religious views on gender roles; after \textit{Mahmoud}, strict scrutiny applies.
\item \textbf{Opposite-Side (No Exemption): Secular Commune.} A Thoreau-inspired commune's claim fails because ``purely secular considerations'' do not trigger the Religion Clauses.
\item \textbf{Opposite-Side (No Exemption): No Schooling at All.} A faith forbidding all learning leaves children illiterate; as White noted, that is ``a very different case.''
\item \textbf{Fence-Sitter: Vaccine Mandate.} A school measles-vaccine rule allows medical but not religious exemptions. \textit{Prince} supports the State, but \textit{Fulton} may make the rule not generally applicable.
\end{enumerate}

\section{Critique}

\textbf{``Good citizen'' bias and the Thoreau boundary.} Relying on the Amish's 300-year pedigree and low crime rate risks disfavoring newer or unpopular faiths and entangles courts in judging religious lifestyles. Excluding secular conscience sharpens the tension between free-exercise accommodation and Establishment Clause neutrality.

\textbf{Children's autonomy and scope.} Deferring children's independent interests, as Douglas warned, and detaching \textit{Yoder} from a sect's existential survival in \textit{Mahmoud} raise acute limiting-principle questions for public-school curricula.

\section{Key Quotations}

\begin{quote}
``The essence of all that has been said and written on the subject is that only those interests of the highest order and those not otherwise served can overbalance legitimate claims to the free exercise of religion.''
\end{quote}
\noindent --- Chief Justice Burger, Opinion of the Court, 406 U.S. at 215.

\begin{quote}
``A way of life, however virtuous and admirable, may not be interposed as a barrier to reasonable state regulation of education if it is based on purely secular considerations; to have the protection of the Religion Clauses, the claims must be rooted in religious belief.''
\end{quote}
\noindent --- Chief Justice Burger, Opinion of the Court, 406 U.S. at 215.

\begin{quote}
``A regulation neutral on its face may, in its application, nonetheless offend the constitutional requirement for governmental neutrality if it unduly burdens the free exercise of religion.''
\end{quote}
\noindent --- Chief Justice Burger, Opinion of the Court, 406 U.S. at 220.

\end{document}
